\textbf{Deriving the penalty from restoration.}
The penalty $U$ follows directly from the restoring displacement in Eq.~\eqref{eq:restoring-step}. For $\rho=0$ or $\rho\geq\epsilon$, we can rewrite that displacement as
\begin{equation}
d_i^R
=-\eta_0\bar r_i
-\eta_{\mathrm{anc}}
\left(\|\bar r_i\|_2-\sqrt D\rho\right)_+
\frac{\bar r_i}{\|\bar r_i\|_2},
\end{equation}
with value zero at $\bar r_i=0$. Integrating this radial field gives
\begin{equation}
U(r)=\frac{\eta_0}{2}\|r\|_2^2+
\frac{\eta_{\mathrm{anc}}}{2}
\left(\|r\|_2-\sqrt D\rho\right)_+^2.
\label{eq:radial-potential}
\end{equation}
Differentiating $U$ recovers $d_i^R=-\nabla_r U(\bar r_i)$. Thus, CAST's restoration is a descent step on a penalty that weakly discourages all departure from the base and adds resistance beyond radius $\rho$.

\textbf{Connecting target regression to regularized critic optimization.}
Using this derived penalty, define
\begin{equation}
\mathcal L_{\mathrm{reg}}(\theta)
=\mathbb E_s\!\left[
\frac{1}{K}\sum_{i=1}^K
\left(
-\frac{\eta_Q\bar Q_\phi(s,z_i,a_\theta(s,z_i))}
{q_{\mathrm{scale}}}
+U(r_\theta(s,z_i))
\right)
\right].
\end{equation}
This loss favors higher predicted return while penalizing departure. With both caps inactive, CAST's target correction is
$\Delta_i=\eta_Q\nabla_a\bar Q_\phi/q_{\mathrm{scale}}-\nabla_r U(\bar r_i)$,
which is the negative action gradient of the corresponding term in $\mathcal L_{\mathrm{reg}}$.
Let $J_i=\partial r_\theta(s,z_i)/\partial\theta$ at the target-construction parameters $\theta_{\mathrm{old}}$. Holding samples, critic parameters, normalization, and regression targets fixed gives
\begin{equation}
\begin{aligned}
\left.\nabla_\theta\mathcal L_{\mathrm{actor}}\right|_{\theta_{\mathrm{old}}}
&=-\frac{2}{KD}\mathbb E_s\!\left[\sum_i J_i^\top\Delta_i\right]\\
&=\frac{2}{D}
\left.\nabla_\theta\mathcal L_{\mathrm{reg}}\right|_{\theta_{\mathrm{old}}}.
\end{aligned}
\label{eq:local-gradient}
\end{equation}
Therefore, unclipped target regression and direct optimization of $\mathcal L_{\mathrm{reg}}$ have the same actor-gradient direction at the parameters used to construct the targets. This establishes how the action corrections implement the balance between critic guidance and preservation. CAST's clipping operations additionally modify each correction before it contributes to the actor gradient.

\textbf{Connecting restoration to distributional departure.}
Restoration needs a measure of departure from the pretrained policy that can be computed without policy likelihoods. The paired residual provides such a measure. At a fixed observation $s$, let $P_\theta^s$ and $P_0^s$ denote the candidate distributions before selection, with finite second moments. Using the same noise pairs each refined action with its base action, with squared distance $\|r_\theta(s,z)\|_2^2$. Since the squared 2-Wasserstein distance minimizes expected squared distance over all pairings, and $U(r)\geq\eta_0\|r\|_2^2/2$, for $\eta_0>0$ we obtain
\begin{equation}
\begin{aligned}
W_2^2(P_\theta^s,P_0^s)
&\leq \mathbb E_z\|r_\theta(s,z)\|_2^2\\
&\leq \frac{2}{\eta_0}\mathbb E_z U(r_\theta(s,z)).
\end{aligned}
\label{eq:coupling-bound}
\end{equation}
Thus, a small expected restoration penalty implies a small distance between the candidate distributions. This gives $U$ a concrete role in the regularized objective: critic guidance favors higher predicted return, while $U$ penalizes an upper bound on departure from the pretrained distribution. Because this departure is measured relative to the frozen base, it captures accumulated change across updates. The component study in Sec.~\ref{sec:exp-controls} provides complementary evidence for this design: clipping alone yields zero success on the five tested LIBERO tasks, while restoration retains nonzero success on every task.
